% =====================================================================
%  Section 8 -- Conclusions
% =====================================================================

\section{Conclusions}
\label{sec:conclusion}

We presented FloodLite, the first systematic CPU-and-browser edge-deployment benchmark for lightweight flood segmentation. Across three lightweight UNet students, four knowledge-distillation configurations, three cross-validation folds, two execution providers, and FP32 $+$ INT8 precisions, we have established four results that we believe will be useful to the operational disaster-response community:

\begin{enumerate}
\item ImageNet-pretrained lightweight UNets match or exceed a UNet$+$EfficientNet-B0 teacher on FSSD \emph{without} distillation, and combined response$+$feature KD does not improve over no-KD on any of the three architectures.
\item Post-training INT8 quantization stability is sharply architecture-dependent, with bounded-activation backbones (ReLU6) quantizing cleanly while HardSwish-based and attention-based backbones require quantization-aware training.
% M1 (pre-submission review): the former "15.8x speedup over the teacher"
% compared a PyTorch-eager teacher (274 ms) against ORT-CPU students. On a common
% common ORT-CPU backend the teacher is 37.6 ms (cross-checked in lat_m2_ort117.json;
% 36.7 ms when re-timed under ORT 1.22.1 -- version-insensitive), so the honest
% INT8-vs-teacher speedup is 2.2x. Numbers below are final; the M2 CPU rows are
% ORT 1.22.1 (see runs/3fold/lat_m2.json _provenance), the browser rows ORT-Web 1.17.3.
\item The recommended deployable configuration, EfficientNet-Lite0-UNet INT8 served by ONNX Runtime CPU, runs at \deployableMTwoFPS~FPS on a 2023 laptop CPU at $3.83\times$ size reduction with IoU 0.885 (95\% of the teacher's accuracy) --- a $2.2\times$ speedup over the teacher on the same ONNX Runtime CPU backend; in the browser, where INT8 is unsupported by ONNX Runtime Web~1.17, an FP32 student runs at 6--7.5\,FPS in a single-threaded WebAssembly sandbox (MobileNetV3-Small fastest).
\item Zero-shot evaluation bounds the deployment claims twice over. Across sensors, every FSSD-trained model collapses on Sentinel-2 optical imagery (IoU $\approx 0.16$). Within the close-range RGB regime, on overlap-verified external data, FP32 students transfer as well as the teacher (0.69--0.72 IoU), but flood-only training data produces models that false-alarm on 72--100\% of water-free scenes, and INT8 degradation grows under domain shift, up to outright collapse --- so in-domain PTQ validation alone cannot qualify a quantized artifact for deployment.
\end{enumerate}

The same \texttt{.onnx} artifact serves both targets, which we offer as a practical pattern for flood-segmentation deployment on heterogeneous edge hardware.

We release our trained checkpoints, the ONNX FP32 and INT8 exports, all benchmarking harnesses, and a Zenodo-archived snapshot of the code repository under permissive licences. Future work includes retraining on negative-image-bearing combined data (Section~\ref{sec:methods-data}) to remove the false-alarm failure mode of Section~\ref{sec:results-zeroshot}, domain adaptation for cross-sensor transfer --- our zero-shot Sen1Floods11 evaluation collapses to IoU $\approx 0.16$ (Section~\ref{sec:results-ood}), so the recommendations hold only for in-domain close-range RGB --- ARM CPU benchmarks on Raspberry~Pi~4 and AWS Graviton, energy measurements, cost-weighted (recall-favouring) evaluation matched to the operational error asymmetry (Section~\ref{sec:discussion-asymmetry}), and quantization-aware-training of the HardSwish and attention-based backbones whose post-training quantization stability we have characterised as a current practitioner trap.
